\section{The bounded removal and veto prescription}
\label{sec:rejection}

This section first states the bounded rule when the reconstructed hard
density is the generated source density. It applies to ordinary matching
and to the global-sum campaign. The sampled Delta rule reconstructs the
\emph{undamped} local hard density instead, with the source factor
retained in the amount removed. Its different physical transfer is
given in Sec.~\ref{sec:delta-sampled}; the common implementation of the
probabilities must not obscure that distinction.

\subsection{The three densities and the eligible shower measure}
For an ordinary source, define on a common labelled real phase space
\begin{equation}
 H(q)=R(q)-D_{\rm eff}(q),\qquad
 H_-(q)=\pos{-H(q)}=\pos{D_{\rm eff}(q)-R(q)}.
 \label{eq:hard-negative}
\end{equation}
Here $D_{\rm eff}=R-H$ is the effective source transfer defined by the
generated hard density. Its name does not assert that it is a single
nonnegative shower kernel. Delta conditioning will introduce additional
labels without changing this principle.

The scalar leading kernel, in the same measure, has the structure
\begin{equation}
 K_g(q)=\sum_h \frac{B_h(b_h)}{J_h}\,
             \kappa_h^{(1)}(r_h)\,g_h(q).
 \label{eq:leading-kernel}
\end{equation}
The sum includes the appropriate Born-flavour and leading-colour
probabilities, luminosity/flux factors, recoil alternatives and the
distribution of allowed starting scales. The indicator or marginal
factor $g_h$ restricts the histories to the profiled matching support.
It includes the implementation's beam bounds, cutoffs and protected
heavy-flavour regions. It is not the source G function. No physical
Sudakov factor is inserted into Eq.~\eqref{eq:leading-kernel}.

For example, backward ISR uses the ratio of momentum-weighted PDFs,
schematically
\begin{equation}
 \dd\kappa_h\ \sim\ \frac{\dd t}{t}\,
 \frac{\as}{2\pi}P(z)\,
 \frac{(x/z)f_a(x/z,\mu)}{xf_b(x,\mu)}\,
 \dd z\,\frac{\dd\phi}{2\pi},\qquad t=\pt^2.
 \label{eq:isr-ratio}
\end{equation}
The displayed $1/z$ cannot be discarded by replacing $xf$ with $f$.
Equation~\eqref{eq:isr-ratio} is only schematic: the actual map, Jacobian,
mass terms, support and scale convention must come from the selected
kernel. In particular, the evaluator's leading coefficient at the
fixed-order scales is not the complete running-$t$, Sudakov-weighted
native branching law. A random starting scale contributes a marginal
such as $\Pr(Q>\sqrt t)$; observing one source event with $Q>\sqrt t$
does not make that marginal identically one.

The strict API requires nonnegative $R$ and $K_g$. The supported
signed-residual mode also admits finite signed PDF luminosities and uses
only $K_g^+=\pos{K_g}$ as transferable support. For $K_g\leq0$ it keeps
both S and any negative-H residual: rejecting a negative leading S
measure cannot compensate negative-H removal. It never replaces the
physical densities by absolute values. A negative nominal S weight
from virtual corrections alone does not determine the sign of $K_g$.
Native positive-part, flavour-threshold and small-denominator
policies also enter the relation between an analytic leading kernel and
the actual accepted native law. Their equivalence must be checked, rather
than assumed from a common grid name.

\subsection{Removing only what a bounded leading veto can transfer}
The desired full H removal would ask for the S veto probability
$H_-/K_g$. To keep a physical probability without silently deleting an
uncompensated part, define
\begin{equation}
 C(q)=\min\{H_-(q),K_g^+(q)\},\qquad
 H_{-,\rm rem}(q)=H_-(q)-C(q).
 \label{eq:bounded-transfer}
\end{equation}
The acceptance in this source-density construction is
\begin{equation}
 a_S(q)=1-\frac{C(q)}{K_g(q)}
       =1-\min\!\left(1,\frac{H_-(q)}{K_g(q)}\right),
 \quad K_g(q)>0.
 \label{eq:s-acceptance}
\end{equation}
For a negative H event the acceptance is instead
\begin{equation}
 a_H^-(q)=\frac{H_-(q)-C(q)}{H_-(q)}
 =\max\!\left(1-\frac{K_g^+(q)}{H_-(q)},0\right),
 \quad H_-(q)>0.
 \label{eq:h-acceptance}
\end{equation}
Positive H events survive. S events with no emission or with a protected
first emission survive. If $K_g\leq0$ and $H_->0$, the signed-residual rule
retains the negative H contribution and performs no S rejection. A
negative H event whose reconstructed density in its qualified partition has no negative part
is an inconsistency and causes a failure, rather than a convenient
change of its classification.

These are whole-event decisions. Rejecting a first branching and
continuing from the original Born state would change the generated
no-emission and subsequent-emission populations and is a different
algorithm. The present procedure does not repeat an event until it
passes, replace its weight by its absolute value, or rescale the retained
sample to the input rate.

If $D_{\rm eff}=K_g$ and the same support applies, then
$a_S=\min(1,R/K_g)$, and Eq.~\eqref{eq:h-acceptance} removes all negative
H events. This is the familiar ideal multiplicative--accumulative limit.
The source and kernel in this implementation are more general, so that
identity is a special case, not the definition imposed on the input.

\subsection{Why the requested ratio can exceed one}
The subtraction was designed to cancel singular expansions, not to be a
pointwise upper bound on a separately profiled scalar shower kernel.
There is no general inequality $H_-\leq K_g$. G/spin completion, the
source low-scale prescription, different support, source luminosities,
and conditional Delta histories can all matter. If a phase-space region
has no eligible shower image, then $K_g=0$ while its source hard density
may still be negative.

As a local arithmetic example, take $R=1$, $D_{\rm eff}=4$ and $K_g=2$
in a common density unit. Full H removal asks for a veto probability
$3/2$. Clipping that probability to one while deleting all negative H
events would discard the missing transfer. Equations
\eqref{eq:bounded-transfer}--\eqref{eq:h-acceptance} transfer two units
and retain one negative hard-density unit. This example illustrates the
bound; it does not specify the actual physical S emission density.

A strict-positive-H mode may stop when the bound fails. The campaign
instead uses the signed-residual mode, which makes the missing transfer
explicit in the accepted sample. Remaining negative H events are thus
part of the declared prescription, not events accidentally bypassing the
filter. Conversely, their small number in one process does not prove
that the finite-coupling compensation is accurate in that process.

\subsection{The exact finite-first-step difference}
\label{sec:finite-difference}
Let $\zeta$ collect the Born kinematics and all shower-relevant source
labels: flavours, colour flow, starting scales and, when applicable,
auxiliary source variables. The signed S source measure is
$\dd\mu_S(\zeta)$. If $T_1(q\mid\zeta)$ is the actual native first
emission law and $g(\zeta,q)$ marks eligible emissions, define
\begin{equation}
 E_g(q)=\int\dd\mu_S(\zeta)\,
                  T_1(q\mid\zeta)\,g(\zeta,q).
 \label{eq:actual-s-density}
\end{equation}
This is a signed cross-section density. It contains the NLO S weights,
the physical no-harder-emission probability, the distribution of starting
states, the actual PDF/coupling choices and every native sampling policy.
It is not automatically $K_g(q)$.

For the ordinary or source-density construction, removing negative H adds $C(q)$ to the
signed endpoint density, while S rejection subtracts
$C(q)E_g(q)/K_g(q)$. Thus
\begin{equation}
 \delta\rho_1(q)=C(q)\left[1-\frac{E_g(q)}{K_g(q)}\right],
 \qquad K_g(q)>0.
 \label{eq:finite-defect}
\end{equation}
At $K_g\leq0$ the transferred amount is zero and no division is used. With
auxiliary conditioning the corresponding expression must be integrated
over the \emph{same} joint conditional measure; one cannot replace
conditional quantities by separately averaged factors. The Delta section
specifies that extension.

More generally, let $T(q,a)$ be the positive signed-measure increment
from removing negative H events, and let $v(q,a)$ reject eligible S
events on the same retained state. The master identity is
\begin{equation}
 \delta\rho_1(q,a)=T(q,a)-v(q,a)E_g(q,a).
 \label{eq:master-defect}
\end{equation}
The scalar rule above has $T=C$ and $v=C/K_g$. In sampled Delta mode,
$T=\overline\Delta_a C_{0,a}$ and $v=C_{0,a}/K_a$, with
$C_{0,a}=\min\{[-H_{0,a}]_+,K_a\}$. Thus its finite defect is different
even though the same bounded probability routine is used. Unobserved
auxiliaries must be integrated with the corresponding source law;
Sec.~\ref{sec:delta-sampled} makes that average explicit.

Equation~\eqref{eq:finite-defect} explains two logically independent
limitations. The bounded transfer controls the part that cannot be
removed with the chosen leading kernel; even where that bound permits
full H removal, the actual S rejection need not compensate it exactly.
For instance, with $H_-=1$, $K_g=4$ and $E_g=3$, full H removal adds one
unit while the veto removes $3/4$ of a unit. The result differs by
$1/4$, despite a perfectly valid probability $1/4$.

Negative S weights are especially relevant. The same probability thins
positive and negative S events in the current implementation. Rejecting
a negative S event raises, rather than lowers, the signed rate. It is
therefore the signed $E_g$, not the number of rejected S events or their
absolute-weight density, that enters Eq.~\eqref{eq:finite-defect}.

If the generated first emission has the required leading correspondence,
then at fixed resolved kinematics
\begin{equation}
 E_g=K_g+\mathcal O(\as^{p+2}),\qquad
 H_-,K_g=\mathcal O(\as^{p+1}),
 \label{eq:leading-correspondence}
\end{equation}
where the Born process starts at $\as^p$. The change in
Eq.~\eqref{eq:finite-defect} is then one order beyond NLO. This is a
conditional fixed-order argument. Large logarithms at small transverse
momentum do not make its expansion uniform, and a large finite numerical
effect is not ruled out. More importantly, a mismatch of the leading
accepted native law and the analytic kernel, for example from an
unqualified map or clipped trial acceptance, does not acquire a
higher-order label simply because it is observed after showering. The
completed histograms are evidence for the specified implementation;
they are not a proof of Eq.~\eqref{eq:leading-correspondence} everywhere.

\subsection{Why multiplying by one Sudakov is insufficient}
An actual emission law has a conditional no-harder-emission factor,
schematically $\Delta_{\rm PS}(t,Q;\zeta)\kappa_h$. After projecting
to the same $q$, different histories have different Born points,
different $\overline B/B$ ratios, different start scales and potentially
different colour states. In general
\begin{equation}
 \int\dd\zeta\,\overline B(\zeta)\,
 \Delta_{\rm PS}(t,Q;\zeta)\kappa(\zeta,q)
 \ne \Delta_{\rm common}(q)\,K_g(q).
 \label{eq:no-common-sudakov}
\end{equation}
Even the average of exponentials over starting states is not the
exponential of an averaged kernel. Adding the source bogus factor to a
denominator, exponentiating its G completion, or inserting the Delta
source factor as if it were the actual shower Sudakov therefore does
not construct the exact finite S density.

The artificial source factor, the physical shower no-emission
probability and the Delta source modification have three separate jobs.
They can have related limiting behaviour, but their equality is not an
implementation identity.

\subsection{What would make the cancellation exact?}
\label{sec:exact-cancellation}
If the actual signed eligible density $E_g(q)$ were known, a common
probability acting on both signs could use
\begin{equation}
 0\leq C_{\rm exact}(q)\leq\min\{H_-(q),E_g(q)\},\qquad
 v_{\rm exact}(q)=\frac{C_{\rm exact}(q)}{E_g(q)},\quad E_g>0.
 \label{eq:exact-veto}
\end{equation}
It would retain the corresponding hard residual
$H_--C_{\rm exact}$. On $E_g\leq0$ this sign-independent rejection
cannot provide a positive compensating transfer; it must set the
transfer to zero or change the construction. Full H removal would require
$E_g\geq H_-$ on its entire negative support.

An additional option is to reject only positive-weight eligible S
events. Let $E_g^+$ denote their positive cross-section density. Then
\begin{equation}
 0\leq C_+(q)\leq\min\{H_-(q),E_g^+(q)\},\qquad
 v_+(q)=C_+(q)/E_g^+(q)
 \label{eq:positive-s-veto}
\end{equation}
would cancel the same removed H measure exactly at this endpoint, while
retaining negative S events. This is an alternative mathematical
construction, not an option implemented in the present campaign.

Both proposals require information absent from the current scalar
denominator: the signed or positive-weight NLO S population, the actual
conditional first-emission law including Sudakovs and scale mixtures,
all inverse histories and their Jacobians, and sufficient common
support. An empirically fitted density ratio would introduce estimation
and finite-sample errors; it would not establish exact equality by
definition. Exact cancellation at finite statistics would additionally
require a coupling or correlated event construction, not two independent
Bernoulli counts expected merely to agree on average.

Exact equality of first-emission momentum distributions is also weaker
than equality after further showering. A sufficient construction matches
the transfer on the full endpoint state $z$, including colours, masses,
beams and imported scale limits, or on equivalence classes with identical
continuation laws. More generally, the necessary condition is that the
signed transfer have vanishing pushforward through the relevant shower;
different endpoint labels can sometimes induce the same continuation.
Momentum-only matching does not establish that condition. An LHE restart
need not retain every internal state variable of an uninterrupted shower.
These issues are not resolved by changing one numerator or denominator in
Eq.~\eqref{eq:s-acceptance}. For the sampled prescription, use the
actual source negative measure on its component space in
Eqs.~\eqref{eq:exact-veto}--\eqref{eq:positive-s-veto}, including the
source Delta average when its auxiliary is marginalized.

\subsection{A useful full-shower decomposition}
Let $I$ be the source measure after its first-step proposal map, $M$ the
same measure after decisions, and $U_{\rm rest}$ the supported restart
continuation. With $U_{\rm direct}$ denoting the uninterrupted reference
from the original source $F_0$, the measured difference is
\begin{equation}
 U_{\rm rest}M-U_{\rm direct}F_0
 =U_{\rm rest}(M-I)
  +\left[U_{\rm rest}I-U_{\rm direct}F_0\right].
 \label{eq:full-decomposition}
\end{equation}
The first term is the removal/veto effect under a common continuation.
The second is measured by the unrejected blue restart control and
includes presentation, first-step profile and restart differences. It
does not isolate restart alone. In particular, the matching profile uses
global-recoil mode 0 for its first step, whereas the direct full reference
uses mode 2 with its standard first-global-branch restrictions.
Eliminating the first term alone would therefore not remove a nonzero
black--blue control difference.
